\section{Excluding residual five and the uniform three-unit margin}
\label{sec:residual-six-margin}

\begin{theorem}\label{thm:residual-five-excluded}
Let $D$ be a finite nonempty strongly connected all-deficient oriented
graph minimal under deletion of arbitrary nonempty arc sets. Then
\[
 R=2M-n=\sum_x\eta(x)\ge6,
 \qquad
 {M\ge\left\lceil\frac n2\right\rceil+3.}
\]
At odd order $R\ge7$ by parity; at even order $R\ge6$.
Equivalently,
$|A(D)|\le\lfloor n(n-2)/2\rfloor-3$.
\end{theorem}
\begin{proof}
The earlier parity-sensitive margin gives $R\ge5$.
Every positive residual support has at least three vertices by
Corollary~\ref{cor:three-positive-residual-support}.
If $R=5$, its only possible positive integer partitions are
\[
 (3,1,1),\quad(2,2,1),\quad(2,1,1,1),\quad(1,1,1,1,1).
\]
They are excluded, respectively, by
Theorem~\ref{thm:three-positive-two-units},
Theorem~\ref{thm:no-two-two-one},
Theorem~\ref{thm:no-two-one-one-one}, and
Theorem~\ref{thm:no-five-unit-residuals}.
Hence $R\ne5$ and $R\ge6$.
Its parity agrees with that of $n$. Therefore integer rounding gives
$M\ge\lceil(n+6)/2\rceil=\lceil n/2\rceil+3$ in both cases.
\end{proof}

\begin{corollary}[Extension to all arc-set-minimal counterexamples]
\label{cor:residual-six-sink}
Let $D$ be any finite nonempty all-deficient oriented graph minimal
under deletion of arbitrary nonempty arc sets. If a sink strongly
connected component has $k$ vertices and $o=n-k$ vertices lie outside
it, then
\[
 M(D)\ge\left\lceil\frac{k}{2}\right\rceil+3
       +\left\lceil\frac{o(k-2)}2\right\rceil
 \ge\left\lceil\frac n2\right\rceil+o+3.
\]
In particular the uniform three-unit margin holds without strong
connectivity, and equality forces strong connectivity.
\end{corollary}
\begin{proof}
The sink component inherits deficiency and arc-set minimality.
Apply Theorem~\ref{thm:residual-five-excluded} internally, and add the
established outside-component estimate
$M(D[V\setminus S])+M(S,V\setminus S)\ge\lceil o(k-2)/2\rceil$.
Since every all-deficient graph has minimum outdegree at least two,
$k\ge5$. Thus the outside term is at least $\lceil3o/2\rceil$.
The old rounding inequality
\[
 \left\lceil\frac{k}{2}\right\rceil
 +\left\lceil\frac{3o}{2}\right\rceil
 \ge\left\lceil\frac{k+o}{2}\right\rceil+o
\]
proves the second bound. If $o>0$ it is strictly stronger than the
uniform bound, giving the equality conclusion.
\end{proof}

At both $n=2\delta+3$ and $n=2\delta+4$, the new missing-pair bound
gives $M\ge\delta+5$.
The next unresolved residual values are six at even order and seven
at odd order. The full conjecture is not settled by this result.
The bound concerns the stated deletion-minimal class; it is not
transferred backwards to an arbitrary denser counterexample.
